%==============================================================
%  book.tex
%  hasklib — mathematical companion / derivations, combined edition
%
%  Single-document version of ito_foundations.tex, gbm.tex, ou.tex,
%  abm.tex, cir.tex, schemes.tex, as chapters in reading order:
%  foundations -> GBM -> OU -> ABM -> CIR -> numerical schemes.
%
%  Skeleton only: chapter shells + shared preamble are wired up;
%  content to be filled in per chapter.
%==============================================================
\documentclass[11pt]{report}

%---------------- packages ----------------
\usepackage[utf8]{inputenc}
\usepackage[T1]{fontenc}
\usepackage{amsmath, amssymb, amsthm}
\usepackage{mathtools}
\usepackage[a4paper, margin=1in]{geometry}
\usepackage{xcolor}
\usepackage{hyperref}
\usepackage{cleveref}
\usepackage{bm}
\usepackage{harvard}   % Harvard author-date referencing (with agsm.bst)

%---------------- notation macros ----------------
% (identical to the notation blocks in ito_foundations.tex / gbm.tex / ou.tex / abm.tex / cir.tex / schemes.tex)
\newcommand{\Prob}{\mathbb{P}}
\newcommand{\Qmeas}{\mathbb{Q}}
\newcommand{\E}{\mathbb{E}}
\newcommand{\Var}{\operatorname{Var}}
\newcommand{\Cov}{\operatorname{Cov}}
\newcommand{\filt}{\mathcal{F}}
\newcommand{\Om}{\Omega}

\newcommand{\dd}{\mathrm{d}}
\newcommand{\dW}{\dd W_t}
\newcommand{\dt}{\dd t}
\newcommand{\R}{\mathbb{R}}

\newcommand{\drift}{a}
\newcommand{\diff}{b}

%---------------- theorem environments ----------------
% Union of every environment used across the six source files,
% numbered within chapter (Theorem 2.1, Definition 3.2, ...).
\theoremstyle{plain}
\newtheorem{theorem}{Theorem}[chapter]
\newtheorem{lemma}{Lemma}[chapter]
\newtheorem{proposition}{Proposition}[chapter]
\newtheorem{corollary}{Corollary}[chapter]

\theoremstyle{definition}
\newtheorem{definition}{Definition}[chapter]
\newtheorem{assumption}{Assumption}[chapter]

\theoremstyle{remark}
\newtheorem{remark}{Remark}[chapter]

\numberwithin{equation}{chapter}

\newcommand{\TODO}[1]{\medskip\noindent\textbf{[\textcolor{red}{TO DERIVE}: #1]}\medskip}

%==============================================================
\title{hasklib: Stochastic Process Derivations}
\author{Hubert Kasinski}
\date{\today}
%==============================================================

\begin{document}
\maketitle

\begin{abstract}
This document collects the stochastic-calculus results that the
stochastic layer of \texttt{hasklib} relies on: from the
probability-space setup, Brownian motion, It\^o SDE, to Geometric Brownian Motion
, Ornstein--Uhlenbeck, Arithmetic Brownian Motion, Cox--Ingersoll--Ross, and numerical schemes.
Standard measure-theoretic probability is assumed; see
\citeasnoun{shreve2004} or \citeasnoun{williams1991} for construction.
\end{abstract}

\tableofcontents



%--------------------------------------------------------------
\chapter{Foundations: It\^o Calculus and the SDE Framework}
%--------------------------------------------------------------
% Source: ito_foundations.tex
\section{Probability space and Brownian motion}

We work on a filtered probability space
$(\Om, \filt, \{\filt_t\}_{t\ge 0}, \Prob)$ satisfying the usual
conditions. The filtration $\{\filt_t\}$ represents the information
available up to time $t$.

\begin{definition}[Brownian motion]
\leavevmode\\A standard Brownian motion $\{W_t\}_{t\ge 0}$ is a process with
\begin{enumerate}
  \item $W_0 = 0$ almost surely;
  \item Independent increments: for $s < t < u $, $W_u - W_t$ is
        independent of $W_t - W_s$;
  \item Gaussian increments: $W_t - W_s \sim \mathcal{N}(0,\, t-s)$;
  \item Continuous sample paths.
\end{enumerate}
\end{definition}

\begin{remark}
       \leavevmode\\
In code, property (3) is what \texttt{NormalRng::draw()} supplies:
a step of length $\Delta t$ uses $W_{t+\Delta t}-W_t = \sqrt{\Delta t}\,Z$
with $Z\sim\mathcal N(0,1)$. This substitution is the single bridge
between the continuous theory here and the discrete simulation code.

\end{remark}
\leavevmode\\
One might notice (2) and (4) pull in
different directions. Independent increments mean each step of the
path is a fresh draw, carrying no memory of the past --- nothing
about $\filt_t$ constrains $W_{t+h}-W_t$. Continuity, on the other
hand, means the path never jumps: $W_{t+h}$ must stay close to $W_t$
as $h\to 0$. It is not obvious a single process can do both at once,
and the price of doing so turns out to be that the path is extremely
rough and nowhere differentiable.
\leavevmode\\
\leavevmode\\
Fixing $t$ and considering the difference quotient over a step $h>0$. By
the Gaussian-increment property, $W_{t+h}-W_t \sim \mathcal N(0,h)$,
so the difference quotient
\[
   \frac{W_{t+h}-W_t}{h}
\]
is itself Gaussian with mean $0$ and variance
\[
   \Var\!\left(\frac{W_{t+h}-W_t}{h}\right)
   = \frac{1}{h^2}\Var(W_{t+h}-W_t)
   = \frac{h}{h^2} = \frac{1}{h},
\]
using that dividing a random variable by a constant $h$ scales its
variance by $1/h^2$. Hence the standard deviation of the difference
quotient is $1/\sqrt h$, which diverges as $h\downarrow 0$: the path
moves by order $\sqrt h$ over a step of order $h$, so no finite slope
survives the limit. This informal argument is made rigorous by the
Paley--Wiener--Zygmund theorem \cite{pwz1933}: with probability one,
$t\mapsto W_t$ is nowhere differentiable. It also follows that the
path has infinite first-order (total) variation on every interval, so
integrals $\int f\,\dd W_s$ cannot be defined pathwise as
Riemann--Stieltjes integrals. Its quadratic variation, however, is
finite, and it is precisely this finite second-order variation that
the It\^o integral and It\^o's lemma are built on.

%--------------------------------------------------------------
\section{Quadratic variation}\label{sec:qv}
%--------------------------------------------------------------

The property that distinguishes stochastic from ordinary calculus is
that Brownian motion accumulates quadratic variation at a
nonzero rate.

\begin{proposition}[Quadratic variation of Brownian motion]\label{prop:qv}
\leavevmode\\
       For a standard Brownian motion, the quadratic variation over $[0,T]$
satisfies
\[
   \langle W\rangle_T \;=\; \lim_{\|\Pi\|\to 0}\sum_{k} \big(W_{t_{k+1}}-W_{t_k}\big)^2
        \;=\; T
   \qquad \text{in } L^2,
\]
where $\Pi$ is a partition of $[0,T]$. Informally this is written
\[
   (\dW)^2 = \dt.
\]
\end{proposition}

\begin{proof}
       \leavevmode\\
Write the partition as $\Pi=\{0=t_0<t_1<\dots<t_n=T\}$ and abbreviate
\[
  \Delta t_k = t_{k+1}-t_k, \qquad
  \Delta W_k = W_{t_{k+1}}-W_{t_k}, \qquad
  S_\Pi = \sum_{k=0}^{n-1}\big(\Delta W_k\big)^2 .
\]
``Convergence in $L^2$'' means the mean-square error vanishes:
$\E\big[(S_\Pi-T)^2\big]\to0$ as the mesh $\|\Pi\|:=\max_k\Delta t_k\to0$.
We use only two facts about the increments, both directly from the definition
of Brownian motion:
\begin{itemize}
  \item Each $\Delta W_k\sim\mathcal N(0,\Delta t_k)$.
  \item Increments over disjoint intervals are independent.
\end{itemize}
\medskip
\noindent We then calculate the mean of $S_\Pi$, i.e $\E[S_\Pi]$
\\
Since $\E[\Delta W_k]=0$, the second moment is the variance:
\[
  \E\big[(\Delta W_k)^2\big]=\Var(\Delta W_k)=\Delta t_k .
\]
Summing over $k$,
\[
  \E[S_\Pi]=\sum_{k}\E\big[(\Delta W_k)^2\big]
           =\sum_k \Delta t_k = T .
\]
This holds for every partition, not merely in the limit. Because the
mean is exactly $T$, the mean-square error is the variance:
\[
  \E\big[(S_\Pi-T)^2\big]=\Var(S_\Pi) .
\]
The whole statement then reduces to showing $\Var(S_\Pi)\to0$.

\medskip
\noindent We then calculate the variance of a single term.
\\
Normalising the increment as $\Delta W_k=\sqrt{\Delta t_k}\,Z$ with
$Z\sim\mathcal N(0,1)$, so $(\Delta W_k)^2=\Delta t_k\,Z^2$ and
\[
  \Var\big((\Delta W_k)^2\big)=(\Delta t_k)^2\,\Var(Z^2).
\]
For a standard normal, $\E[Z^2]=1$ and $\E[Z^4]=3$, hence
$\Var(Z^2)=\E[Z^4]-\big(\E[Z^2]\big)^2=3-1=2$. (The fourth moment is one
Gaussian integration by parts: with the density $\phi$ satisfying
$\phi'(z)=-z\,\phi(z)$, integrating $\int z^4\phi\,\dd z$ by parts gives
$\E[Z^4]=3\,\E[Z^2]=3$.) Therefore,
\[
  \Var\big((\Delta W_k)^2\big)=2\,(\Delta t_k)^2 .
\]

\medskip
\noindent\emph{Step 3}\; Variance of the sum, then let the mesh shrink.
\\
The increments $\Delta W_k$ are independent, meaning the squared increments
$(\Delta W_k)^2$ are independent too, and the variance of a sum of
independent variables is the sum of the variances:
\[
  \Var(S_\Pi)=\sum_k \Var\big((\Delta W_k)^2\big)=2\sum_k (\Delta t_k)^2 .
\]
Bounding one factor of $\Delta t_k$ in each term by the mesh:
\[
  \sum_k (\Delta t_k)^2 \le \Big(\max_k \Delta t_k\Big)\sum_k \Delta t_k
                        = \|\Pi\|\,T .
\]
Combining,
\[
  \E\big[(S_\Pi-T)^2\big]=\Var(S_\Pi)\le 2\,\|\Pi\|\,T
  \;\xrightarrow[\;\|\Pi\|\to0\;]{}\;0 ,
\]
which is exactly the claimed $L^2$ convergence.
\end{proof}

\begin{remark}
\leavevmode\\
The mean computation ($\E[S_\Pi]=T$) shows the average comes out
exactly right; the variance computation shows it comes out reliably,
not just on average. This second part matters because a single term
is not close to its own mean: $(\Delta W_k)^2$ has standard deviation
$\sqrt2\,\Delta t_k$, about the same size as its mean $\Delta t_k$
itself. The reason the full sum $S_\Pi$ still ends up close to $T$ is
that the terms are independent, so their ups and downs cancel out
when added together --- this cancellation is what the variance
computation shows, and it is a property of the sum, not of any single
term. Inside a sum (equivalently, an integral), the random
$(\Delta W_k)^2$ may therefore be replaced by the deterministic
$\Delta t_k$: this replacement is the entire content of the shorthand
$(\dW)^2=\dt$ used in the next section.
\end{remark}

%--------------------------------------------------------------
\section{The It\^o stochastic differential equation}
%--------------------------------------------------------------

\begin{definition}[It\^o SDE]
\leavevmode\\
An It\^o process $X_t$ solves the stochastic differential equation
\begin{equation}\label{eq:sde}
   \dd X_t \;=\; \drift(t,X_t)\,\dt \;+\; \diff(t,X_t)\,\dW,
\end{equation}
interpreted as the integral equation
\[
   X_t = X_0 + \int_0^t \drift(s,X_s)\,\dd s
              + \int_0^t \diff(s,X_s)\,\dd W_s,
\]
where the second integral is an It\^o integral. We call $\drift$ the
drift and $\diff$ the diffusion coefficients.
\end{definition}

\begin{remark}[Correspondence with the code]
  \leavevmode\\
Equation~\eqref{eq:sde} is exactly what the \texttt{StochasticProcess}
base class encodes: \texttt{drift(t,x)} returns $\drift(t,x)$ and
\texttt{diffusion(t,x)} returns $\diff(t,x)$. Every subsequent
process (GBM, OU, CIR) is a choice of these two functions.
\end{remark}

\begin{assumption}[Existence and uniqueness]\label{ass:existence-uniqueness}
  \leavevmode\\
We assume $\drift$ and $\diff$ satisfy the standard Lipschitz and
linear-growth conditions guaranteeing a unique strong solution to
\eqref{eq:sde}; see \citeasnoun{oksendal2003}, Thm.~5.2.1. 
\end{assumption}

%--------------------------------------------------------------
\section{It\^o's lemma}
%--------------------------------------------------------------

It\^o's lemma is the chain rule for functions of an It\^o process, and
it is the engine behind every closed-form solution used in
\texttt{hasklib} (the GBM exact sampler, the OU solution).

\begin{theorem}[It\^o's lemma, one dimension]\label{thm:ito}
\leavevmode\\
Let $X_t$ solve \eqref{eq:sde} and let $\varphi(t,x)\in C^{1,2}$. Then
\begin{equation}\label{eq:ito}
   \dd \varphi(t,X_t) =
     \Big( \frac{\partial\varphi}{\partial t} + \drift\, \frac{\partial\varphi}{\partial x} + \frac{1}{2} \diff^2 \frac{\partial^2\varphi}{\partial x^2} \Big)\dt
     \;+\; \diff\, \frac{\partial\varphi}{\partial x} \,\dW,
\end{equation}
\end{theorem}



\begin{proof}[Informal derivation]
  \leavevmode\\
This is the standard second-order Taylor argument. It is not a proof in the
strict sense, as a rigorous version requires the construction of the It\^o
integral \cite[Sec.~4.2.1]{shreve2004}, which is outside the scope of this document.
The derivation stems around the idea that, because $\dd X$ carries a Brownian term of size
$\sqrt{\dt}$, a contribution that ordinary calculus discards as ``second
order'' is in fact first order and should be kept.

\medskip
\noindent
Over a step of length $\dt$ the increment $\dW$ has typical size $\sqrt{\dt}$
(it is $\mathcal N(0,\dt)$).
\\Ranking the elementary products by their power of
$\dt$:
\[
  \dt = O(\dt), \qquad
  \dW = O(\dt^{1/2}), \qquad
  (\dW)^2 = O(\dt), \qquad
  \dt\,\dW = O(\dt^{3/2}), \qquad
  (\dt)^2 = O(\dt^2).
\]
Discarding everything of order smaller than $\dt$, and using \cref{prop:qv} to
replace $(\dW)^2$ by its in-sum limit $\dt$, gives the It\^o
multiplication table
\begin{equation}\label{eq:itotable}
  (\dW)^2=\dt, \qquad \dt\,\dW=0, \qquad (\dt)^2=0 .
\end{equation}
The last two entries are not literally zero; they are negligible next to $\dt$
and vanish in the mesh limit for exactly the reason the fluctuations did in
\cref{sec:qv}.

\medskip
\noindent Taylor expansion:
\leavevmode\\
Expand $\varphi(t+\dt,\,X_t+\dd X)$ about $(t,X_t)$ to second order:
\[
  \dd \varphi = \frac{\partial\varphi}{\partial t}\,\dt + \frac{\partial\varphi}{\partial x}\,\dd X
        + \frac{1}{2} \frac{\partial^2\varphi}{\partial t^2}\,(\dt)^2
        + \frac{\partial^2\varphi}{\partial t\partial x}\,\dt\,\dd X
        + \frac{1}{2} \frac{\partial^2\varphi}{\partial x^2}\,(\dd X)^2
        + \cdots
\]
Substitute the SDE $\dd X = \drift\,\dt + \diff\,\dW$ and square:
\[
  (\dd X)^2 = \drift^2(\dt)^2 + 2\,\drift\diff\,\dt\,\dW + \diff^2(\dW)^2 .
\]
Apply \eqref{eq:itotable}. In $(\dd X)^2$ the first two terms die and only
$\diff^2(\dW)^2=\diff^2\,\dt$ survives:
\[
  (\dd X)^2 = \diff^2\,\dt .
\]
The two remaining second-order terms of the expansion vanish for the same
reason:
\[
  \frac{1}{2} \frac{\partial^2\varphi}{\partial t^2}\,(\dt)^2 = 0, \qquad
  \frac{\partial^2\varphi}{\partial t\partial x}\,\dt\,\dd X
     = \frac{\partial^2\varphi}{\partial t\partial x}\big(\drift\,(\dt)^2 + \diff\,\dt\,\dW\big) = 0 .
\]
Leaving:
\[
  \dd \varphi = \frac{\partial\varphi}{\partial t}\,\dt + \frac{\partial\varphi}{\partial x}\,\dd X + \frac{1}{2} \frac{\partial^2\varphi}{\partial x^2}\,\diff^2\,\dt .
\]
Substitute $\dd X = \drift\,\dt + \diff\,\dW$ once more and collect the $\dt$
and $\dW$ parts:
\[
  \dd \varphi = \Big( \frac{\partial\varphi}{\partial t} + \drift\, \frac{\partial\varphi}{\partial x} + \frac{1}{2} \diff^2 \frac{\partial^2\varphi}{\partial x^2} \Big)\dt
        + \diff\, \frac{\partial\varphi}{\partial x} \,\dW ,
\]
which is \eqref{eq:ito}.
\end{proof}

\begin{remark}
\leavevmode\\
Had $X$ been a smooth (finite-variation) path, $(\dd X)^2$ would have been
genuinely $O((\dt)^2)$ and dropped out, leaving the ordinary chain rule
$\dd \varphi = \frac{\partial\varphi}{\partial t}\,\dt + \frac{\partial\varphi}{\partial x}\,\dd X$. The whole gap between stochastic and ordinary
calculus is the single surviving term $\frac{1}{2}\diff^2 \frac{\partial^2\varphi}{\partial x^2}\,\dt$, and it
survives only because $(\dW)^2=\dt$ rather than $0$ --- i.e.\ because Brownian
motion has nonzero quadratic variation. Every closed-form result in
the code that ``corrects a drift by half a variance''
is exactly this one term.
\end{remark}

%--------------------------------------------------------------
\chapter{Geometric Brownian Motion}
%--------------------------------------------------------------
% Source: gbm.tex

%--------------------------------------------------------------
\section{The GBM SDE}
%--------------------------------------------------------------

\begin{definition}[Geometric Brownian motion]\label{def:gbm}
\leavevmode\\
  $X_t$ follows geometric Brownian motion with drift $\mu\in\R$ and
volatility $\sigma>0$ if it solves
\begin{equation}\label{eq:gbmsde}
   \dd X_t = \mu X_t\,\dt + \sigma X_t\,\dW, \qquad X_0>0.
\end{equation}
Referring back to the previous chapter on It\^o Calculus, this is the choice
$\drift(t,x)=\mu x$, \\$\diff(t,x)=\sigma x$ in the general SDE
$\dd X_t = \drift(t,X_t)\dt+\diff(t,X_t)\dW$.
\end{definition}

\begin{remark}[Correspondence with the code]
  \leavevmode\\
This is exactly why the pair of functions \texttt{GBM::drift} and
\texttt{GBM::diffusion} return: \texttt{drift(t,x)} $=\mu x$,
\texttt{diffusion(t,x)} $=\sigma x$. \texttt{EulerMaruyama} consumes
these polymorphically through the \texttt{StochasticProcess}
interface; the sampler derived below is the exact alternative
to that discretisation, used by \texttt{GBM::sample\_terminal} instead
of stepping the Euler scheme forward.
\end{remark}
\leavevmode\\
\noindent Both $\drift$ and $\diff$ are affine in $x$, hence globally
Lipschitz and of linear growth on all of $\R$ (\cref{ass:existence-uniqueness}),
so \eqref{eq:gbmsde} has a unique strong solution for any $X_0\in\R$. In
particular this holds for $X_0>0$; positivity of $X_t$ itself is
established below once the explicit solution is derived.

%--------------------------------------------------------------
\section{Linearising via It\^o's lemma}
%--------------------------------------------------------------

The SDE~\eqref{eq:gbmsde} has no constant-coefficient closed form as
written, because both coefficients scale with $X_t$. The standard fix
is to apply It\^o's lemma to $\varphi(t,x)=\ln x$: this removes the
$x$-dependence from both coefficients and turns \eqref{eq:gbmsde} into
an SDE with constant drift and diffusion, which integrates
directly.

\begin{theorem}[Exact terminal law of GBM]\label{thm:gbm-exact}
\leavevmode\\
Let $X_t$ solve \eqref{eq:gbmsde}. Then for any $T>0$,
\begin{equation}\label{eq:gbmexact}
   X_T = X_0\exp\!\left(\Big(\mu-\tfrac12\sigma^2\Big)T + \sigma W_T\right),
\end{equation}
where $W_T\sim\mathcal N(0,T)$.
\end{theorem}

\begin{proof}
\leavevmode\\
Apply It\^o's lemma (\cref{thm:ito}) to
$\varphi(t,x)=\ln x$,
so:
\\$\frac{\partial\varphi}{\partial t}=0$, $\frac{\partial\varphi}{\partial x}=\frac{1}{x}$, $\frac{\partial^2\varphi}{\partial x^2}=-\frac{1}{x^2}$. With
$\drift(t,X_t)=\mu X_t$ and $\diff(t,X_t)=\sigma X_t$,
\[
   \dd(\ln X_t)
     = \Big( 0 + \mu X_t\cdot\frac1{X_t}
             + \tfrac12\,\sigma^2 X_t^2\cdot\Big(\!-\frac1{X_t^2}\Big) \Big)\dt
       + \sigma X_t\cdot\frac1{X_t}\,\dW .
\]
The $X_t$ factors cancel completely, leaving
\begin{equation}\label{eq:lnlinear}
   \dd(\ln X_t) = \Big(\mu-\tfrac12\sigma^2\Big)\dt + \sigma\,\dW .
\end{equation}
The right-hand side no longer depends on $X_t$: it is Brownian motion
with constant drift $\mu-\tfrac12\sigma^2$ and constant volatility
$\sigma$, so \eqref{eq:lnlinear} integrates directly from $0$ to $T$:
\[
   \ln X_T - \ln X_0
     = \Big(\mu-\tfrac12\sigma^2\Big)T + \sigma\,(W_T-W_0)
     = \Big(\mu-\tfrac12\sigma^2\Big)T + \sigma W_T,
\]
using $W_0=0$. Exponentiating gives \eqref{eq:gbmexact}. Since $W_T\sim
\mathcal N(0,T)$ by definition of Brownian motion, this is a complete
description of the law of $X_T$ given $X_0$.
\end{proof}

\begin{remark}
  \leavevmode\\
Equation~\eqref{eq:gbmexact} is sampled with a single normal draw: generate
$Z\sim\mathcal N(0,1)$, set $W_T=\sqrt T\,Z$, and evaluate the
right-hand side. No time-stepping is needed, which is what makes it
an exact sampler rather than a discretisation, and it is exactly
what \texttt{GBM::sample\_terminal} implements using
\texttt{NormalRng::draw()}.
\end{remark}

\begin{remark}[Where the $-\tfrac12\sigma^2$ comes from]
  \leavevmode\\
The drift that appears in \eqref{eq:gbmexact} is not $\mu$, the
drift of $X_t$ itself, but $\mu-\tfrac12\sigma^2$. That correction is
exactly the $\tfrac12\diff^2 \frac{\partial^2\varphi}{\partial x^2}$ term of It\^o's lemma, which is the term
with no analogue in ordinary calculus, mentioned in the previous chapter. If one (incorrectly) applied the
ordinary chain rule to $\ln X_t$, the correction would be missing and
the resulting sampler would be biased upward in expectation, as
\cref{prop:gbm-mean} below makes precise.
\end{remark}

%--------------------------------------------------------------
\section{Moments}
%--------------------------------------------------------------

The moment-matching test in \texttt{test\_stochastic.cpp} checks
simulated terminal samples against the following closed-form mean and
variance.

\begin{lemma}[Moment generating function of a Gaussian]\label{lem:mgf}
\leavevmode\\
If $Y\sim\mathcal N(m,v)$ then $\E[e^Y]=\exp\!\big(m+\tfrac12 v\big)$.
\end{lemma}


\begin{proof}
\leavevmode\\
Write $Y=m+\sqrt v\,Z$ with $Z\sim\mathcal N(0,1)$. Complete the square
in the Gaussian integral:
\[
   \E[e^Y] = e^m\,\E\big[e^{\sqrt v\,Z}\big]
           = e^m \int_{-\infty}^{\infty}
               \frac{1}{\sqrt{2\pi}}\,e^{\sqrt v\,z-\frac12 z^2}\,\dd z
           = e^m\,e^{\frac12 v}
             \int_{-\infty}^{\infty}
               \frac{1}{\sqrt{2\pi}}\,e^{-\frac12(z-\sqrt v)^2}\,\dd z
           = e^{m+\frac12 v},
\]
the last integral being $1$ since it is a $\mathcal N(\sqrt v,1)$
density integrated over $\R$.
\end{proof}

\begin{proposition}[Mean of GBM]\label{prop:gbm-mean}
\leavevmode\\
$\E[X_T] = X_0\,e^{\mu T}$.
\end{proposition}

\begin{proof}
\leavevmode\\
By \cref{thm:gbm-exact}, $X_T=X_0\,e^Y$ with
$Y=(\mu-\tfrac12\sigma^2)T+\sigma W_T$.
\\
Since $W_T\sim\mathcal N(0,T)$,
$Y\sim\mathcal N\big((\mu-\tfrac12\sigma^2)T,\ \sigma^2 T\big)$.

\noindent Applying
\cref{lem:mgf} with $m=(\mu-\tfrac12\sigma^2)T$ and $v=\sigma^2 T$,
\[
   \E[X_T] = X_0\,\E[e^Y]
           = X_0\exp\!\Big( \big(\mu-\tfrac12\sigma^2\big)T + \tfrac12\sigma^2 T \Big)
           = X_0\,e^{\mu T}.
\]
Where the $-\tfrac12\sigma^2 T$ It\^o correction and the $+\tfrac12\sigma^2 T$
from the moment generating function cancel exactly, leaving the mean
growing at the "natural" rate $\mu$ despite the correction appearing
inside every individual path.
\end{proof}

\begin{proposition}[Variance of GBM]\label{prop:gbm-var}
\leavevmode\\

$\Var(X_T) = X_0^2\,e^{2\mu T}\big(e^{\sigma^2 T}-1\big)$.
\end{proposition}
\medskip
\begin{proof}
\leavevmode\\
$\Var(X_T)=\E[X_T^2]-\big(\E[X_T]\big)^2$.
\\For the second moment, write
$X_T^2=X_0^2\,e^{2Y}$ with $Y$ as above, so $2Y\sim\mathcal
N\big(2(\mu-\tfrac12\sigma^2)T,\ 4\sigma^2 T\big)$. By \cref{lem:mgf}
with $m=2(\mu-\tfrac12\sigma^2)T$ and $v=4\sigma^2 T$,
\[
   \E[X_T^2] = X_0^2\exp\!\Big( 2\big(\mu-\tfrac12\sigma^2\big)T + 2\sigma^2 T \Big)
             = X_0^2\,e^{(2\mu+\sigma^2)T}.
\]
Subtracting $\big(\E[X_T]\big)^2=X_0^2\,e^{2\mu T}$ from
\cref{prop:gbm-mean} and factoring,
\[
   \Var(X_T) = X_0^2\,e^{(2\mu+\sigma^2)T} - X_0^2\,e^{2\mu T}
             = X_0^2\,e^{2\mu T}\big(e^{\sigma^2 T}-1\big). \qedhere
\]
\end{proof}

\begin{remark}[What the test checks]
  \leavevmode\\
In the code, \texttt{test\_stochastic.cpp} draws a large number of terminal samples
from \texttt{GBM::sample\_terminal}, computes their sample mean and
sample variance, and checks these against
\cref{prop:gbm-mean,prop:gbm-var} to within Monte Carlo tolerance. This
is a check on the sampler, not on \eqref{eq:gbmsde} itself ---
it would catch, for instance, a sign error in the $-\tfrac12\sigma^2$
correction, since that error changes the mean but leaves the
qualitative shape of the distribution intact.
\end{remark}

%--------------------------------------------------------------
\chapter{The Ornstein--Uhlenbeck Process}
%--------------------------------------------------------------
% Source: ou.tex

%--------------------------------------------------------------
\chapter{Arithmetic Brownian Motion}
%--------------------------------------------------------------
% Source: abm.tex

%--------------------------------------------------------------
\chapter{The Cox--Ingersoll--Ross Process}
%--------------------------------------------------------------
% Source: cir.tex

%--------------------------------------------------------------
\chapter{Numerical Schemes: Euler--Maruyama and Milstein}
%--------------------------------------------------------------
% Source: schemes.tex

%--------------------------------------------------------------
% References cited across the source files, all already present
% in refs.bib: shreve2004, oksendal2003, kloeden1992, williams1991,
% pwz1933, brigomercurio2006, cir1985, lord2010.
%--------------------------------------------------------------
\bibliographystyle{agsm}
\bibliography{refs}

\end{document}
